% Section 2. Categorical Evidential Bayesian Neural Network
% Backbone: Discussion 6.2 (sec:discussion-cat-ebnn). See STORY.md.

\section{Categorical Evidential Bayesian Neural Network}

% Slide 2.1. Source: Methods 4.3.1 (sec:ebnn-model, eq:ebnn-hierarchy),
% 4.3.2 (sec:ebnn-likelihood-equivalence, eq:ebnn-likelihood-equivalence).
\begin{frame}{The Cat-eBNN Model}
	\textbf{Message.} The Cat-eBNN is the Dirichlet network of EDL with a posterior over the shared network weights.

	\medskip
	\textbf{TODO: content.} Model sketch. Likelihood is the Dirichlet mean $m_y(x;w)$.
\end{frame}

% Slide 2.2. Source: Background 2.9.3 (sec:bg-dirichlet-total-concentration,
% eq:aleatoric-as-function-of-total-concentration, eq:aleatoric-limits). Appendix A.1.
\begin{frame}{The Total Concentration Parameter Sets the Allocation}
	\textbf{Message.} At a fixed mean vector, a larger total concentration parameter gives a higher aleatoric term and a lower distributional term.

	\medskip
	\textbf{TODO: content.} Figure of both terms against $\alpha_0$ at a fixed mean vector. New figure from \texttt{scripts/}.
\end{frame}

% Slide 2.3. Source: Background 2.3 (sec:bg-scale-invariance, eq:scale-invariance).
% Methods 4.3.3 (sec:ebnn-identifiability, eq:ebnn-likelihood-invariance),
% 4.3.4 (sec:gamma-prior, eq:assumed-total-concentration).
\begin{frame}{Single Labels Cannot Inform It. The Gamma Prior}
	\textbf{Message.} The single label likelihood is invariant to the total concentration parameter at a fixed mean vector. So we state the assumption with a Gamma prior.

	\medskip
	\textbf{TODO: content.} Scale invariance in one line. Gamma prior with prior mode and prior standard deviation. Honest scope sentence.
\end{frame}

% Slide 2.4. Source: Experiments 5.1.1 to 5.1.3, 5.1.5 (sec:exp-setup,
% sec:exp-models-training, sec:exp-training-setup, sec:exp-metrics-single).
\begin{frame}{Experimental Setup with Single Labels}
	\textbf{Message.} We compare the Cat-eBNN and the Cat-BNN on Fashion-MNIST and CIFAR-10 with SGLD.

	\medskip
	\textbf{TODO: content.} Datasets. One SGLD chain, fixed budget. Cold, warm, shifted initialisation. Prior sweeps. Metrics.
\end{frame}

% Slide 2.5. Source: Experiments 5.2.1 (sec:exp-starts, tab:e2-start-procedure,
% fig:start-trajectory). Discussion 6.2. Candidate for backup.
\begin{frame}{Initialisation}
	\textbf{Message.} Warm and shifted initialisation give higher validation accuracy than cold initialisation. Shifted initialisation avoids the accuracy loss when the prior mode is far above the checkpoint.

	\medskip
	\textbf{TODO: content.} Trajectory figure or table extract. One line on finite sampling budget.
\end{frame}

% Slide 2.6. Source: Experiments 5.2.2 (sec:exp-prior-mode, tab:e2-mode-sweep),
% 5.2.3 (sec:exp-prior-sd, tab:e2-sd-sweep). Discussion 6.2.
\begin{frame}{Prior Mode and Prior Standard Deviation}
	\textbf{Message.} The realised total concentration follows the prior mode. A narrower prior keeps it closer. A higher realised total concentration gives a higher aleatoric share and a lower distributional share.

	\medskip
	\textbf{TODO: content.} Extracts of the mode sweep and sd sweep tables. One line on the fixed mean effect.
\end{frame}

% Slide 2.7. Source: Experiments 5.2.4 (sec:exp-predictive, tab:e2-predictive),
% 5.2.5 (sec:exp-decomposition, tab:e2-terms). Discussion 6.2.
\begin{frame}{Predictive Performance and the Contribution}
	\textbf{Message.} The Cat-eBNN has competitive accuracy but no consistent predictive improvement over the Cat-BNN. Its contribution is the complete uncertainty representation and the explicit prior assumption.

	\medskip
	\textbf{TODO: content.} Extract of the predictive table. Two closing lines on contribution and scope.
\end{frame}
